\begin{center}
	\large{\textbf{KATA PENGANTAR}}
\end{center}

\indent
Segala puji syukur dipanjatkan kehadirat Allah SWT karena telah melimpahkan rahmat, hidayah, petunjuk, kekuatan, dan rezeki kepada penulis sehingga penulis dapat menyelesaikan Tesis yang berjudul:
\begin{center}
	\large{\textbf{MATRIKS KOEFISIEN BAGI UKURAN KETERBELITAN KEADAAN DUA DAN TIGA QUBIT}}
\end{center}
Tesis ini disusun sebagai salah satu syarat untuk memperoleh gelar Magister (S2) pada Program Studi Fisika di Institut Teknologi Sepuluh Nopember. Penulisan tesis ini tidak lepas dari bantuan, bimbingan, dan dukungan dari berbagai pihak. Oleh karena itu, pada kesempatan ini penulis ingin menyampaikan ucapan terima kasih dan penghargaan yang setinggi-tingginya kepada semua pihak yang telah membantu dan memberi motivasi kepada penulis, baik secara langsung maupun tidak langsung, dalam proses penyusunan tesis ini:
\begin{enumerate}
	
	\item Ibu yang senantiasa memberikan dukungan, doa, dan harapan kepada penulis.
	
	\item Prof. Agus Purwanto, D.Sc., selaku dosen pembimbing atas bimbingan, arahan, masukan, dan motivasinya kepada penulis.
	
	\item Bapak Dr. Lila Yuwana, M.Si. selaku Kepala Departemen Fisika yang telah memberikan kemudahan selama proses kuliah.
	
	\item  Para guru dan dosen yang telah ikhlas mengajar, memberikan ilmu, dan memberikan inspirasi kepada penulis.
	
	\item Teman seperjuangan Laboratorium Fisika Teori dan Filsafat Alam (LaFTiFA) selama penulis menempuh magister di Fisika ITS. 
	
	\item Seluruh pihak yang tidak bisa disebutkan satu-persatu yang selalu memberikan dukungan dan motivasi dalam menyelesaikan tesis ini.
\end{enumerate}

\indent Penulis menyadari bahwa dalam penulisan ini masih terdapat kekurangan, sehingga segala kritik dan saran yang sifatnya membangun sangat diharapkan untuk kesempurnaan Tesis ini.
Penulis berharap Tesis ini dapat bermanfaat bagi penulis sendiri pada khususnya dan pembaca pada umumnya.




\begin{flushright}
	Surabaya, Juli 2026\\
	\vspace{2cm}
	Penulis
\end{flushright}
\newpage

\newpage
\thispagestyle{empty}
\begin{center}
	\topskip0pt
	\vspace*{\fill}
	\textit{\textbf{"Halaman ini sengaja dikosongkan"}}
	\vspace*{\fill}
\end{center}
\pagebreak